\thispagestyle{empty}
\setlength{\parindent}{0in}
\setlength{\parskip}{0em}

\textbf{POLITECHNIKA RZESZOWSKA} \hfill Rzeszów 2027 \\
\textbf{im. Ignacego Łukasiewicza} \\
\textbf{WYDZIAŁ MATEMATYKI I FIZYKI STOSOWANEJ} \\

\begin{center}
	\textbf{STRESZCZENIE  PRACY  DYPLOMOWEJ }
\end{center}

\textbf{Tytuł:} Wpływ promptu na progi bólu agenta LLM \\
\textbf{Autor:} Kamil Piejko \\
\textbf{Promotor:} dr inż. Marcin Kowalik, prof. PRz \\
\textbf{Słowa klucze:} \dower{do 5 słów kluczowych}
\\

\dower{tekst streszczenia w języku polskim, pisany na końcu}


% Wstawienie pionowej przestrzeni do połowy strony
\vfill
$$\noindent\rule{\textwidth}{2pt}$$



\textbf{RZESZOW UNIVERSITY OF TECHNOLOGY} \hfill Rzeszów 2027 \\
\textbf{FACULTY OF MATHEMATICS AND APPLIED PHYSICS } \\

\begin{center}
	\textbf{DIPLOMA THESIS ABSTRACT}
\end{center}

\textbf{Title:} The Influence of Prompt Wording on the Pain Thresholds of an LLM Agent \dower{tytuł angielski: propozycja, do zatwierdzenia}\\
\textbf{Author:} Kamil Piejko \\
\textbf{Supervisor:} Marcin Kowalik, PhD Eng., Assoc. Prof. \dower{zapis stopnia promotora po angielsku do potwierdzenia} \\ %% Angielski zapis stopni/tytułów naukowych promotora 

\textbf{Key words:} \dower{up to 5 key words}
\\

\dower{tekst streszczenia w języku angielskim}

\vfill

% Poniżej przykłady angielskich odpowiedników polskich stopni i tytułów naukowych promotora:
% - doktor | dr → Dr Firstname Lastname lub Firstname Lastname, PhD
% - dr inż. → Dr Eng. Firstname Lastname lub Firstname Lastname, PhD Eng.
% - doktor habilitowany | dr hab. → Firstname Lastname, PhD, DSc
% - dr hab. inż. → Firstname Lastname, PhD, DSc Eng.
% - dr hab. prof. PRz → Firstname Lastname, PhD, DSc, Assoc. Prof.
% - prof. dr hab. → Prof. Firstname Lastname
% - prof. dr hab. inż. → Prof. Eng. Firstname Lastname
% - doktor, profesor PRz → Doctor, Firstname Lastname, Associate Professor lub Firstname Lastname, PhD, Assoc. Prof.
